\section{Implications and Limitations}\label{discussion}
This section discusses implications for student rocketry teams and identifies the limitations of this study.

% \subsection{Interpretation of Results}
% \subsubsection{Non-blocking Orchestrator one-tick delay}
% The apogee overshoot observed in both the Snapshot and Rate Group strategies is a structural property of the non-blocking Orchestrator design. At callback $n$, the Orchestrator sends the current sensor packet to the FSW and immediately returns the actuation value computed at callback $n-1$ from the actuation buffer, without waiting for the FSW to process the packet from callback $n$. The actuation response corresponding to callback $n$ is therefore not applied until callback $n+1$, producing a one-tick delay, where the actuation applied at any callback is always one step behind the sensor reading that triggered it. This delay directly affects the \texttt{AirbrakeController}, as the predicted apogee estimate and the air brake deployment level are always one step behind the true simulation state, causing the PI controller to under-correct during the coasting phase and leading to a systematic overshoot. These results confirm $H_1$. 

% All three strategies remain well within the 12\,m acceptance threshold defined in Section~\ref{ss:rqs}, confirming that each is suitable for pre-flight validation despite their architectural differences. The nature of this delay was further confirmed by running the framework against a second rocket configuration, in which the overshoot magnitude differed, but the delay persisted across both non-blocking strategies, confirming that it is a property of the Orchestrator design rather than the rocket configuration. 

% The timing behaviour produced by this delay is representative of real embedded flight software execution. On real hardware, sensor readings will always be slightly stale by the time they are processed on a fixed hardware clock. Although both strategies share this behaviour, the source of the delay differs. In Snapshot, it arises from the non-blocking Orchestrator, while in Rate Group, it additionally arises from the mismatch between the two independent clocks, which is why Rate Group shows higher variance than Snapshot in both SIL and HIL contexts, most prominently in HIL. Lockstep, on the contrary, represents an idealised synchronisation model that does not exist in real hardware interacting with its environment.

% \subsubsection{Wall-clock time and execution overhead}
% It is evident from Table \ref{tab:summary} and Figure \ref{fig:wall_time_violin} that HIL configurations are slower than their SIL counterparts. For Lockstep, the overhead is most notable in both contexts due to its blocking nature. The additional cost in HIL relative to SIL (48.8\,s versus 44.9\,s) is attributable to the Ethernet round-trip between the host and the Raspberry Pi 4, as well as the reduced CPU performance of the target hardware, extending the processing time of the FSW. For the Snapshot and Rate Group strategies, neither blocks the simulation, so the additional overhead is attributed to just the Ethernet transmission latency, as reflected in the smaller but consistent gap in wall-clock times. These results confirm $H_2$: HIL configurations produce longer wall-clock times than their SIL counterparts, attributable to Ethernet latency and reduced CPU performance on the target hardware.

% \subsubsection{Rate Group Frequency Sweep}
% The Rate Group frequency sweep confirms that apogee error decreases monotonically with tick frequency, from 209\,m at 1\,Hz to 4.1\,m at 50\,Hz. At higher tick frequencies, the FSW reads more recent sensor data and actuates more frequently, reducing the window of stale data between ticks and allowing the \texttt{AirbrakeController} to track the trajectory of the rocket with higher accuracy. The diminishing accuracy gains and extra cost beyond 25\,Hz informed the selection of this frequency for the comparison of coupling strategies in Section~\ref{s:evaluation}.

% \subsubsection{Fault Injection} \label{sec:fault_inj}
% Following fault injection illustrated in Figure \ref{fig:fault_injection}, the \texttt{RecoveryController} detected the barometric freeze at $t=11.5$\,s and switched to IMU-based free-fall detection. The drogue parachute was deployed at $t=23.60$\,s in SIL and $t=23.61$\,s in HIL, approximately $0.1$\,s after apogee at $t=23.5$\,s. These results confirm $H_3$. The FSW successfully detected the sensor fault and deployed the drogue parachute near apogee in both execution contexts. 

% Notably, the IMU-based fallback logic deployed the drogue significantly closer to apogee than the nominal barometric detection. Under no-fault conditions, drogue deployment occurred 1.54\,s after apogee, while the fallback path reduced this gap by 1.44\,s.

% This unintentional finding further demonstrates the value of closed-loop testing in flight software development. This result could not have been observed through open-loop testing, as it depends on the FSW to feed the actuation output back into the simulation state. The slight apogee overshoot of 3069\,m relative to the nominal Snapshot apogee of 3066\,m is a separate consequence of the barometric freeze. With altitude held constant, the \texttt{AirbrakeController} can no longer compute a valid predicted apogee, causing under-deployment of the air brakes during the coasting phase.

% The difference in detection performance arises from the confirmation window of 10 consecutive pressure readings for barometric detection. In contrast, IMU-based detection responds as soon as vertical acceleration drops below \texttt{ACCEL\_FREEFALL\_THRESH}, a condition satisfied within one callback call at apogee. The main parachute was deployed at $t=43.60$\,s in both SIL and HIL, exactly $20$\,s after drogue deployment, through the \texttt{MAIN\_DEPLOY\_TIMEOUT} fallback, confirming that the complete recovery sequence was completed under faulty-sensor conditions.

% Together, these results answer RQ2 and RQ3. Coupling strategy and execution context measurably affect apogee accuracy, deployment timing and wall-clock overhead ($H_1$, $H_2$), yet all three strategies remain within the pre-flight acceptance threshold of 12\,m in both SIL and HIL. The fault injection results confirm $H_3$: the framework successfully injects a representative sensor fault, and the FSW recovers in both execution contexts.

\subsection{Implications for Student Rocketry Teams} \label{implications}
The primary audience for RITL is university rocketry teams operating under the constraints of limited budgets, part-time contributors and single-use rockets, where a software fault during flight cannot be corrected. 



%The result of this study suggests that RITL is a practical tool for this context.

To assess practical accessibility, the framework was informally trialled as part of a recruitment exercise for the RISE student rocketry team, in which prospective members were evaluated for selection in the following year. Participants were provided with an RITL deployment, including a base implementation of the FSW that connects to RocketPy and were tasked with independently implementing parachute deployment and air brake control logic. The trial was conducted in SIL only. Participants successfully executed SIL runs with accurate parachute deployments and control logic, suggesting that the framework is accessible to newcomers without prior knowledge of the internals. A controlled usability study across multiple teams would be needed to substantiate this claim more rigorously.

The choice of coupling strategy should reflect the validation goal. Lockstep is best suited for validating flight logic correctness under ideal conditions, particularly control algorithm behaviour, as its tight synchronisation eliminates actuation delay and produces an idealised closed-loop response similar to model-based validation tools such as Simulink. The advantage over such tools is that the software under test is the actual flight software that will fly, not a model of it. There are no code generation steps, no reimplementation of flight logic from block diagrams and no fidelity gap introduced by translating a validated model into a deployed binary.

Snapshot and Rate Group strategies are more suited for timing validation, as both reproduce the one-tick actuation delay that is analogous to real embedded execution. Snapshot is the simpler of the two and serves as a middle ground. Rate Group goes further by running the FSW on an independent hardware timer, making it most useful for characterising how timing non-determinism affects performance before launch. 

The Rate Group tick frequency is not a parameter unique to RITL. Real flight computers already run their control loop at a fixed frequency, chosen by the team based on what their hardware can sustain alongside its other responsibilities. Teams should configure \texttt{RateGroupDriver} to the frequency they have already selected or plan to select for their real flight computer. This ensures that pre-flight validation reflects the actual rate at which recovery logic will execute in flight.

The same principle extends to the other timing parameters RITL exposes: the controller callback rate, the sensor model sampling rate, the parachute trigger rate and the solver's integration time step. The controller callback rate, the Rate Group tick frequency and the parachute trigger rate conceptually represent the same underlying quantity, the frequency of the flight computer's control loop. Since \texttt{AirbrakeController} and \texttt{RecoveryController} both execute within that same loop on real hardware (Section~\ref{flight-logic}), all three should therefore be matched to whatever loop frequency the team's real flight computer runs at. However, RocketPy registers the air brake and parachute callbacks independently, so RITL does not currently enforce this correspondence. In this study, both were held equal at 10 Hz by configuration choice. The sensor model sampling rate instead corresponds to the physical sensor's own output rate, typically specified in its datasheet, and should be matched to that value. The solver's integration time step has no equivalent on real hardware at all, as the physical rocket's flight is continuous rather than stepped. It should instead be chosen fine enough to keep the underlying physics accurate, independent of the target flight computer.

By design, the Orchestrator architecture also lowers the barrier to HIL testing. Since the SIL-to-HIL transition requires no changes to the Orchestrator or flight software beyond recompiling the F Prime deployment for the target hardware, teams could in principle validate their software on actual flight hardware without building expensive test infrastructure. This claim follows from the framework's architecture described in Section~\ref{closed-loop-sim-env} rather than from empirical trial data, since the informal trial above was conducted in SIL only. 

The fault injection capability adds another dimension that is otherwise difficult to test before launch. By having a dedicated fault injection interface in the Orchestrator, teams can apply any sensor fault pattern to the simulation without modifying the FSW, testing recovery logic against failure modes that would otherwise only be discovered during flight.


Finally, the Orchestrator design decouples flight software from the simulator, allowing either to be replaced independently without modifying the other side of the interface. Teams that switch trajectory simulators, upgrade their FSW framework, or retarget to different hardware need only update the relevant side of the coupling.

These architectural properties and the positive outcome of the SIL trial suggest RITL is a practical tool for student rocketry teams.

\subsection{Limitations} \label{limitations}
The following limitations concern constraints in the current implementation and experimental scope, rather than the validity of the reported results, which are addressed separately in Section~\ref{sec:threats}.


% The communication layer between the Orchestrator and the FSW uses TCP in both SIL and HIL execution contexts. An alternative would be UDP, which removes the connection overhead and more closely represents the non-acknowledged transport behaviour of real embedded systems. This would also likely reduce wall-clock time in both execution contexts.

The fault injection experiment targets only the \texttt{RecoveryController}. The \texttt{AirbrakeController} did not implement fallback logic. When tested with a barometric freeze, it held onto the last predicted apogee estimate, causing under-deployment of the air brakes and the overshoot observed in the fault injection results.

% The RocketPy callback rate is fixed at 10\,Hz across all experiments. This rate determines the precision of the simulation loop and directly affects the one-tick delay observation. The sensitivity of the results to this parameter was not evaluated; thus, running the framework at different callback rates would clarify how this design choice affects coupling accuracy across strategies.

The \texttt{MAIN\_DEPLOY\_DELTA\_HPA} threshold is configured to deploy the main parachute approximately 330\,m below apogee, resulting in the main parachute deployment at approximately 2700\,m AGL across all coupling strategies. This does not reflect a typical competition deployment altitude and was not optimised for a specific use case, as the flight software in this study serves as a validation system for the closed-loop framework rather than a production-ready flight computer.

The PI controller was not formally tuned. The constants were selected to produce stable actuation sufficient to exercise the coupling strategies. Any difference between the achieved apogee and the 3000 m target is consistent across all configurations, as the same controller is used in both the baseline and RITL setups, and therefore does not affect the coupling strategy comparison.

The solver time step sensitivity check in Section \ref{callback-sample-rate} was evaluated with a single run per time step, rather than the 10 runs used in the rest of the study. This was a scope reduction given its exploratory purpose and the considerably higher wall-clock cost of the finest step tested. The results were sufficient to show that apogee was not sensitive to step size within the range tested, but a fuller sweep could characterise this relationship more accurately.

